\documentclass[11pt]{article}

\usepackage[margin=1in]{geometry}
\usepackage{amsmath,amssymb,amsthm,mathtools}
\usepackage[hidelinks]{hyperref}

\newtheorem{theorem}{Theorem}
\newtheorem{lemma}[theorem]{Lemma}
\newtheorem{proposition}[theorem]{Proposition}
\theoremstyle{remark}
\newtheorem{remark}[theorem]{Remark}

\newcommand{\dd}{\,\mathrm d}
\newcommand{\one}{\mathbf 1}

\title{Atlas 43: An Improved Global Graphon Interval\\
via Joint Absorption of the Bad Chord and the $C_5$ Surplus}
\author{}
\date{August 24, 2026}

\begin{document}
\maketitle

\begin{abstract}
Let $H$ be the house graph, equivalently the theta graph with path
lengths $1,2,3$.  The earlier global proof established its claimed
Goodman-style bound for
\[
 p\geq \frac{1+1/\sqrt3}{2}=0.7886751345\ldots.
\]
We improve the global threshold to
\[
 p\geq p_0=0.708512194994\ldots,
\]
where $p_0$ is the unique zero in $(7/10,1)$ of
\[
 5440p^4-10848p^3+7981p^2-2580p+309.
\]
The proof applies to every graphon.  Its new ingredient is a joint estimate:
the exact $C_5$ surplus controls complement-degree variance, while the
two codegree factors in the bad chord term are bounded by the square of the
ordinary degree.  A $2\times2$ positive-semidefinite calculation then
absorbs the whole bad term at the stated quartic threshold.
\end{abstract}

\section{Statement}

Let $W\colon[0,1]^2\to[0,1]$ be a symmetric measurable graphon and put
\[
 p:=t(K_2,W),\qquad q:=1-p,\qquad r:=2p-1.
\]
For the house graph $H$, the conjectured lower bound is
\begin{equation}\label{eq:target}
 t(H,W)\geq B(p),
 \qquad
 B(p):=pr(p^3+q^3).
\end{equation}

Define
\begin{equation}\label{eq:quartic}
 P(p):=5440p^4-10848p^3+7981p^2-2580p+309.
\end{equation}

\begin{theorem}[Improved global interval]\label{thm:main}
Let $p_0$ be the unique root of $P$ in $(7/10,1)$.  Then
\[
 p_0=0.7085121949942608\ldots.
\]
Every graphon $W$ of edge density $p\geq p_0$ satisfies
\[
 \boxed{
 t(H,W)\geq p(2p-1)\bigl(p^3+(1-p)^3\bigr).
 }
\]
In particular, the decimal statement $p\geq0.7085121950$ is safe.
\end{theorem}

The proof is self-contained modulo three inequalities already proved in
\texttt{notes/house\_partial.tex}: the exact two-chord identity, the
three-triangle gluing bound, and the displayed square certificate for
$C_5$.  We restate precisely the parts needed below.

\section{The previous decomposition}

Put $U:=1-W$, and let
\[
 C(x,y):=\int_0^1 W(x,s)W(s,y)\dd s
\]
be the two-step kernel.  Set
\[
 A(x,y):=W(x,y)C(x,y),
 \qquad
 D(x,y):=U(x,y)C(x,y).
\]
The two relevant nonnegative five-vertex integrals are
\begin{align}
 F&:=\int W(y,z)A(x,y)A(x,z)\dd x\dd y\dd z,\label{eq:F}\\
 G&:=\int W(y,z)D(x,y)D(x,z)\dd x\dd y\dd z.\label{eq:G}
\end{align}
The exact two-chord decomposition is
\begin{equation}\label{eq:two-chord}
 2t(H,W)=F+t(C_5,W)-G.
\end{equation}

Write
\[
 \tau:=t(K_3,W),
 \qquad
 \delta:=\tau-pr.
\]
Goodman's inequality gives $\delta\geq0$.  The three-triangle gluing
inequality gives $F\geq\tau^3/p^2$, and therefore
\begin{equation}\label{eq:F-surplus}
 F-pr^3
 \geq
 3r^2\delta+\frac{3r}{p}\delta^2+\frac1{p^2}\delta^3.
\end{equation}
The target has the exact split
\begin{equation}\label{eq:target-split}
 2B(p)=pr^3+(p^5-pq^4).
\end{equation}
Consequently
\begin{equation}\label{eq:gap-start}
 \begin{aligned}
 2\bigl(t(H,W)-B(p)\bigr)
 \geq{}&
 3r^2\delta+\frac{3r}{p}\delta^2+\frac1{p^2}\delta^3\\
 &+\bigl(t(C_5,W)-(p^5-pq^4)\bigr)-G.
 \end{aligned}
\end{equation}
The earlier proof used only $G\leq\delta$ and discarded the $C_5$
surplus.  The following joint estimate is the improvement.

\section{Complement variance and transitivity defect}

Let
\[
 d_U(x):=\int_0^1U(x,y)\dd y,
 \qquad
 g(x):=d_U(x)-q,
 \qquad
 v:=\int_0^1g(x)^2\dd x.
\]
Thus $v$ is the variance of the complement degree.  Split it as
\[
 v=v_-+v_+,
 \qquad
 v_-:=\int_{\{g<0\}}g^2,
 \qquad
 v_+:=\int_{\{g\geq0\}}g^2.
\]
Also put
\begin{align}
 n_x&:=\int_{[0,1]^2}
 U(x,y)U(x,z)W(y,z)\dd y\dd z,\label{eq:nx}\\
 n&:=\int_0^1n_x\dd x.\label{eq:n}
\end{align}
Both $v,n$ are nonnegative.  Expanding the triangle density of
$W=1-U$ gives the exact Goodman-surplus identity
\begin{equation}\label{eq:delta-vn}
 \boxed{\delta=2v+n.}
\end{equation}

The sharp $C_5$ proof in \texttt{notes/house\_partial.tex} writes
$T=T_U$, $T\one=q\one+g$, and completes a square.  Keeping the
degree-variance term rather than discarding the whole surplus yields the
following bound.

\begin{lemma}[Retained $C_5$ variance]\label{lem:c5-variance}
For $p\geq1/2$,
\begin{equation}\label{eq:c5-variance}
 t(C_5,W)-(p^5-pq^4)\geq c(p)v,
\end{equation}
where
\begin{equation}\label{eq:c-def}
 c(p):=8q^2-10q+\frac{55}{16}
      =8p^2-6p+\frac{23}{16}.
\end{equation}
\end{lemma}

\begin{proof}
The exact certificate retained in the proof of the sharp $C_5$ bound is
\[
 4\left\|Tg+\frac{8q-5}{8}g\right\|_2^2
 +\left(8q^2-10q+\frac{55}{16}\right)\|g\|_2^2.
\]
The only preceding relaxation replaces the complement $5$-cycle density
by a $4$-path density, and its omitted remainder is nonnegative.  Dropping
the displayed square and that remainder proves
\eqref{eq:c5-variance}.
\end{proof}

We next improve the pointwise estimate $C\leq1$ used in the earlier
bound on $G$.

\begin{lemma}[Degree-sensitive bad-chord bound]\label{lem:G-bound}
With the notation above,
\begin{equation}\label{eq:G-bound}
 G\leq p^2n+2pq\sqrt{v_-n}+q^2v_-+v_+.
\end{equation}
\end{lemma}

\begin{proof}
The two-step kernel satisfies, pointwise,
\[
 C(x,y)
 =\int W(x,s)W(s,y)\dd s
 \leq\int W(x,s)\dd s
 =d_W(x)=p-g(x).
\]
Therefore $D(x,y)=U(x,y)C(x,y)\leq U(x,y)(p-g(x))$, and
\[
 G\leq\int_0^1(p-g(x))^2n_x\dd x.
\]
The negative contribution to $\int g n_x$ is the only one that can
increase the last integral.  On $\{g<0\}$ we have
\[
 d_U(x)=q+g(x)\leq q,
 \qquad
 n_x\leq d_U(x)^2\leq q^2.
\]
Consequently Cauchy--Schwarz gives
\[
 -\int g(x)n_x\dd x
 \leq\int_{\{g<0\}}|g(x)|n_x\dd x
 \leq q\sqrt{v_-n}.
\]
The same pointwise bounds give
\[
 \int g(x)^2n_x\dd x
 \leq q^2v_-+v_+.
\]
Expanding $(p-g)^2$ proves \eqref{eq:G-bound}.
\end{proof}

Combining Lemmas~\ref{lem:c5-variance} and \ref{lem:G-bound},
\begin{equation}\label{eq:joint-bound}
 G-\bigl(t(C_5,W)-(p^5-pq^4)\bigr)
 \leq
 p^2n+2pq\sqrt{v_-n}
 +(q^2-c(p))v_-+(1-c(p))v_+.
\end{equation}

\section{The $2\times2$ absorption and the quartic}

Set
\[
 L:=3r^2=3(2p-1)^2.
\]
In view of $\delta=n+2v_-+2v_+$, first consider the variables
$(n,v_-)$.  Their contribution in \eqref{eq:joint-bound} is at most
$L(n+2v_-)$ precisely when
\begin{equation}\label{eq:quadratic-psd}
 (L-p^2)n-2pq\sqrt{nv}
 +\bigl(2L-(q^2-c(p))\bigr)v\geq0,
\end{equation}
where in this display $v=v_-$.  The remaining $v_+$ contribution is
at most $2Lv_+$ whenever
\begin{equation}\label{eq:positive-variance}
 2L-(1-c(p))\geq0.
\end{equation}
This is a quadratic form in $(\sqrt n,\sqrt v)$.  Its two diagonal
entries are
\begin{align}
 L-p^2&=11p^2-12p+3,\label{eq:diag1}\\
 2L-(q^2-c(p))&=\frac{496p^2-448p+103}{16}.\label{eq:diag2}
\end{align}
Its determinant is
\begin{equation}\label{eq:det}
 \begin{aligned}
 &(L-p^2)\bigl(2L-(q^2-c(p))\bigr)-p^2q^2\\
 &\hspace{35mm}=\frac1{16}P(p).
 \end{aligned}
\end{equation}

For completeness, $P(7/10)=-103/100$ and $P(1)=302$.  Moreover
\[
 P''(p)=65280p^2-65088p+15962
\]
is positive and increasing on $[7/10,1]$, and
$P'(7/10)=2763/25>0$.  Thus $P$ is strictly increasing on this interval
and has a unique root $p_0\in(7/10,1)$.  Direct rational evaluations give
\[
 P(0.7085121949)<0,
 \qquad
 P(0.7085121950)>0,
\]
and hence
\[
 0.7085121949<p_0<0.7085121950.
\]
All polynomial identities and rational endpoint signs in this section are
reproduced by
\texttt{experiments/atlas43\_global\_interval\_verify.py}.

For $p\geq p_0>7/10$, both \eqref{eq:diag1} and
\eqref{eq:diag2} are positive, while \eqref{eq:det} is nonnegative.
The quadratic form \eqref{eq:quadratic-psd} is therefore positive
semidefinite.  Also \eqref{eq:positive-variance} holds: its left-hand
side is $(512p^2-480p+103)/16$, whose larger zero is
$(15+\sqrt{19})/32<7/10$.  We have proved
\begin{equation}\label{eq:absorption}
 G-\bigl(t(C_5,W)-(p^5-pq^4)\bigr)
 \leq3r^2\delta.
\end{equation}

\begin{proof}[Proof of Theorem~\ref{thm:main}]
Substitute \eqref{eq:absorption} into \eqref{eq:gap-start}.  The linear
terms cancel, leaving
\[
 2\bigl(t(H,W)-B(p)\bigr)
 \geq
 \frac{3r}{p}\delta^2+\frac1{p^2}\delta^3\geq0.
\]
Here $p\geq p_0>1/2$, $r\geq0$, and $\delta\geq0$.  This proves the
claimed graphon inequality.
\end{proof}

\section{What changed and what remains}

The old estimate treated the two terms separately:
\[
 t(C_5,W)-(p^5-pq^4)\geq0,
 \qquad G\leq\delta.
\]
It therefore required $3r^2\geq1$.  The new proof retains the exact
degree-variance payment inside the $C_5$ certificate and replaces
$C\leq1$ by $C(x,y)\leq d_W(x)$.  The resulting joint cost is at most
$3r^2\delta$ already at the smaller quartic threshold.

This theorem is global: there is no step-graphon, regularity, finite-rank,
balanced-join, or bounded-perturbation assumption.  It enlarges the proved
portion of the required interval $[1/2,1]$ from
\[
 2\left(1-\frac{1+1/\sqrt3}{2}\right)
 =42.2649\ldots\%
\]
to
\[
 2(1-p_0)=58.29756100\ldots\%.
\]
The interval
\[
 \frac12\leq p<p_0
\]
remains open for arbitrary graphons.  The balanced-join low-side theorem in
\texttt{notes/atlas43\_balanced\_join\_low\_side.tex} remains complementary:
it covers a much larger low-density interval, but only under its structural
complete-join hypothesis.

\end{document}
